\scanpage{11}
\diagram{\begin{array}{ccccc}G''\supset H''&Q&H'\subset G'&P&G_k\end{array}}
\[
\left.\begin{aligned}
P&=\underline{\operatorname{Isom}}(T_B\otimes_{\unsure{Q}}k,T_{DR})\\
Q\subset P'&=\underline{\operatorname{Isom}}(T_{DR},T_{Hdg})\simeq Q\mathbin{\times^{H'}}G'\\
P''&=\underline{\operatorname{Isom}}(T_{Hdg},T_B\otimes k)
\end{aligned}\right\}\text{déf/$k$}
\qquad
\boxed{\begin{gathered}\alpha\in P(C)\\\beta\in Q(C)\subset P'(C)\\\gamma\in P''(C)\end{gathered}}
\qquad
\begin{gathered}\beta\alpha=\gamma^{-1}\\\gamma\beta=\alpha^{-1}\\\alpha\gamma=\beta^{-1}\end{gathered}
\]
\diagram{\begin{tikzcd}
&T_{Hdg}\otimes C\arrow[dr,"\gamma"]&\\
T_{DR}\otimes_k C\arrow[ur,"\beta"]&&T_B\otimes_Q C\arrow[ll,"\alpha"]
\end{tikzcd}}
\[
\left.\begin{aligned}
\mathrm{Fil}&=\underline{\mathrm{Filt}}_{\substack{\otimes,\deleted{\unsure{f}}\text{dans}\\\unsure{\text{classe d'int.}}\text{ de}\\\text{la filtration}\\\text{de }T_{DR}}}\ T_{\unsure{b|B}}\\
\mathrm{Big}&=\underline{\mathrm{Bigr}}_{\substack{\otimes,\text{dans}\\\unsure{\text{classe int.}}\text{ de}\\\text{la }\deleted{\unsure{\text{filtration}}}\inserted{\unsure{\text{bigraduation}}}\\\text{de }T_{DR}}}\ T_B
\end{aligned}\right\}\text{déf/$Q$}
\]
\margin{--- NB Soit $F\supset\mathrm{Fil},\Phi\supset\mathrm{Big}$ définis comme ci-contre, mais sans condition de forme intégrale. Alors $G$ y opère, et $A=F/G,\unsure{B=\Phi/G|\Phi/G}$ s.t des schémas \uline{\unsure{étales}} sur $Q$. Alors $\varphi\in F(C),\psi\in\Phi(C)$ \deleted{définissent} des \deleted{\unsure{$V=$}} $\dot\varphi\in A(C),\dot\psi\in B(C)$. \deleted{\unsure{K}} Soient $\varphi_1\in A,\psi_1\in B$ leurs images dans $A$ et $B$ (pts fermés, \unsure{réduits}). Par définition, $\mathrm{Fil}=F_{\varphi_1},\mathrm{Big}=\Phi_{\psi_1}$.}
\margin{\diagram{\begin{array}{ccc}U'&\subset B'&\subset G'\\\text{rad.}&\text{parabolique}&\\\text{unipotent}&\text{associé à}&\\\text{de }\unsure{U'|B'}&\text{filtration de}&\\&T_{DR}&\end{array}}
$Q=$ fibré des sous groupes de Lévi de $B'$.
\diagram{\begin{tikzcd}B'/U'\simeq L'&G_m\arrow[l,"j"]\end{tikzcd}}
(hom central (Tate))}
\margin{\uline{NB} Si on a $\varepsilon:\deleted{\unsure{\mathrm{det}}}G'\longrightarrow G_m$ on a $(\varepsilon|(B'/U'))\circ j=\mathrm{id}$. donc $j$ \unsure{m.m.} Si $H'\subset B'$ est l'image inverse de $\operatorname{Im}\unsure{j|j'}$ \inserted{dans} $B'$, alors \deleted{\unsure{le tord}} $\unsure{j|j'}$ est connu quand on connaît $H'$.}
\margin{en fait, $\dot\varphi,\dot\psi$ sont dans $A(k),B(k)$, donc dans $A(k_1),B(k_1)$, où $k_1$ est la clôture alg. de $Q$ dans $k$. D'ailleurs, $F\longrightarrow\Phi$ induit $A\overset{\sim}{\longrightarrow}B$, donc}
\margin{constr. explicite en termes du groupe \unsure{déployé} associé au \unsure{réseau cor.|réseau cov.} $\Gamma$ avec $W$ (Weyl) agissant: on tord $\Gamma/W$ par \unsure{autom.} extérieurs....}

\[
P\times P'\longrightarrow P'',\quad P'\times P''\longrightarrow P,\quad P''\times P\longrightarrow P',\qquad Q\subset P'.
\]
\diagram{\begin{tikzcd}\mathrm{Big}\arrow[r,"p"]&\mathrm{Fil}\end{tikzcd}}
\deleted{\unsure{a}}
fibration, faisant de $B$ un torseur sur $\mathrm{Fil}$, de \unsure{gr.pe} $U_{\mathrm{Fil}}$, un groupe unipotent affine sur $F$, de dim relative égale à $\dim U'$. $\mathrm{Fil}$ est de dim égale à $\dim G'/B'$, $\mathrm{Big}$ est de dim égale à $\dim G'/B'+\dim U'=\dim G'/L'$, i.e. $L'=\operatorname{Centr}(j'_0)$, $j'_0:G_m\underset{k}{\longrightarrow}H'$ relevant $j'$.

On a un diagramme commutatif
\diagram{\begin{tikzcd}
P''\arrow[d,"\deleted{\unsure{q}}r"']&P\arrow[d,"q"]\\
\mathrm{Big}_k\arrow[r,"p_k"']&\mathrm{Fil}_k
\end{tikzcd}}
\[
\boxed{\begin{gathered}\varphi=q(\alpha)\in\mathrm{Fil}(C)\\\psi=r(\gamma)\in\mathrm{Big}(C)\end{gathered}}
\qquad [\varphi=p(\psi)]
\]
\margin{$\dim G=g$\\$\dim U'=u$.\\$\dim B'=b$\\$\dim L'=\ell$\\$(u+\ell=b)$}
\diagram{\begin{tikzcd}[column sep=small,row sep=large]
&k(\alpha,\beta)=k(\beta,\gamma)=k(\gamma,\alpha)=K\arrow[dl,dash,"\leq u"']\arrow[d,dash,"\leq u"]\arrow[dr,dash,"\leq g"]&&\\
k(\gamma)\arrow[d,dash,"\leq\ell"']&k(\alpha)\arrow[d,dash,"\leq b"]&k(\beta)\arrow[d,dash,"\leq u"]&\\
k(\psi)\arrow[r,dash,"\leq u"]\arrow[d,dash,"\mathrm{comp}"']&k(\varphi)\arrow[r,dash,"\leq g-b"]\arrow[d,dash,"\mathrm{comp}"]&k\arrow[dd,dash]\arrow[dr,dash]&\\
Q(\psi)\arrow[r,dash,"\leq u"']&Q(\varphi)\arrow[dr,dash,"\leq g-b"]& &\deleted{\unsure{Q(\varphi)}}\ \substack{Q(\varphi_1)\\=Q(\psi_1)}\arrow[dl,dash,"\mathrm{alg}"']\\
&\deleted{\unsure{Q}}&Q&
\end{tikzcd}}
\[
\deleted{\unsure{k(\alpha,\beta)}}K\subset\deleted{\unsure{k(\alpha)}}Q\bigl(k(\alpha),\overline{k(\alpha)}\bigr)
\]
\[
(2i\pi)^n\in k(\alpha)\quad\text{si }Q(-n)\in\operatorname{Ob}M
\]
\margin{$Q(\varphi_1)$ est \uline{indépendant} du plongement de $k$ dans $C$, et \unsure{donc} \illegible{3}, \unsure{définit} \illegible{4}}

\uline{Conj. de transcendance forte}

Si $k\subset\overline Q$, alors ? $\operatorname{degr.\ tr.}\underbrace{k(\alpha,\beta)}_K/k=g+u$
\[
\Longrightarrow\left\{\begin{aligned}
&\operatorname{deg\ tr.}\ k(\alpha)/\unsure{K|k}=g\\
&\operatorname{deg.\ tr.}\ k(\beta)/\unsure{K|k}=u\\
&\operatorname{degr.}\ k(\gamma)/k=g\\
&\operatorname{deg.\ tr.}\ K/k(\beta)=g,\quad\unsure{\text{et}}\\
&\operatorname{deg\ tr.}\ K/k(\alpha)=\operatorname{deg\ tr}\ K/k(\gamma)=u\\
&\operatorname{deg\ tr}\ Q(\varphi)/Q=g-b\\
&\operatorname{deg\ tr.}\ Q(\psi)/Q=g-b+u=g-\ell
\end{aligned}\right.
\]
\mspage{10}
